\documentclass{article}

% Настройки русского языка и шрифтов
\usepackage[utf8]{inputenc}
\usepackage{graphicx}
\usepackage{hyperref}
\usepackage[T2A]{fontenc}
\usepackage[russian]{babel}

\title{Мини-конспект по теме: Теорема Пифагора}
\author{Я Сам}
\date{16 июля 2025 г.}

\begin{document}

\maketitle
\tableofcontents
\section{ВВедение}
Теорема Пифагора — одна из важнейших теорем евклидовой геометрии. Она находит применение в самых разных областях:
\begin{itemize}
    \item  геометрия и тригонометрия
    \item физика
    \item инженерные расчёты
    \item компьютерная графика
\end{itemize}
\section{Формулировка теоремы}
\textbf{Слова:} В прямоугольном треугольнике квадрат гипотенузы равен сумме квадратов катетов.
\begin{equation}
\label{eq:pyfa}
c^{2} = a^{2} + b^{2}
\end{equation}
Как видно из формулы \ref{eq:pyfa}, знание двух сторон позволяет найти третью.

\section{Доказательство (набросок)}
Одно из доказательств основывается на площади квадрата, составленного из четырёх одинаковых прямоугольных треугольников и малого квадрата в центре. Раскладывая площадь двумя способами, получаем $c^{2} = a^{2} + b^{2}$

\section{Примеры расчёта}
\textbf{Пример 1}
\[
a = 3, b = 4
\]
\[
c = \sqrt{a^{2} + b^{2}} = \sqrt{9 + 16} = 5
\]

\textbf{Пример 2}
\begin{enumerate}
    \item Дано: $a = 5$, $b = 12$
    \item Решение: 
    \[
    c = \sqrt{5^{2} + 12^{2}} = \sqrt{25 + 144} = 13
    \]
\section{Таблица значений}
\begin{center}
\begin{tabular}{|c|c|c|}
\hline
Катет $a$ & Катет $b$ & Гипотенуза $c$ \\
\hline
3 & 4 & 5 \\
5 & 12 & 13 \\
7 & 24 & 25 \\
\hline
\end{tabular}
\end{center}


\section{Иллюстрация}
Ниже пример изображения (не забудьте добавить файл \texttt{triangle.png} в ту же папку):
\begin{center}
    \includegraphics[width=0.5\textwidth]{triangle.png}
\end{center}
\section{Заключение}
Теорема Пифагора — один из краеугольных камней геометрии, помогающий решать множество практических задач.

\section{Ссылки и литература}
\begin{itemize}
    \item \href{https://ru.wikipedia.org/wiki/%D0%A2%D0%B5%D0%BE%D1%80%D0%B5%D0%BC%D0%B0_%D0%9F%D0%B8%D1%84%D0%B0%D0%B3%D0%BE%D1%80%D0%B0}{википидия: теорема пифагора}
    \item классические учебники геометрии
\end{itemize}

\end{enumerate}
\end{document}


